\section*{Capítulo \ref{ch:g3:balanced-diet} --- Una alimentación equilibrada}

\begin{solution}{exo:g3:balanced-diet:1}
Aportar \omterm{prop:g3:balanced-diet:jobs}{energía} y aportar material de construcción. \omterm{prop:g3:balanced-diet:jobs}{Energía}: cereales
y féculas (o grasas). Construcción: carne, pescado y \omterm{def:g2:animals-grow-up:born}{huevos}, o
lácteos.
\end{solution}

\begin{solution}{exo:g3:balanced-diet:2}
Fruta y verdura; cereales y féculas; lácteos; carne, pescado y \omterm{def:g2:animals-grow-up:born}{huevos};
grasas; alimentos dulces (la familia de los caprichos). La única
bebida necesaria es el agua.
\end{solution}

\begin{solution}{exo:g3:balanced-diet:3}
Arroz, pan: cereales y féculas. Yogur: lácteos. Zanahorias, manzanas:
fruta y verdura. \omterm{def:g2:animals-grow-up:born}{Huevos}: carne-pescado-huevos. Mantequilla: grasas.
\end{solution}

\begin{solution}{exo:g3:balanced-diet:4}
Porque el cuerpo no para nunca: el corazón late toda la noche, el
pecho sigue respirando y el cuerpo se mantiene caliente a unos
$\qty{37}{\degreeCelsius}$ --- todo pagado en \omterm{prop:g3:balanced-diet:jobs}{energía}.
\end{solution}

\begin{solution}{exo:g3:balanced-diet:5}
Las judías verdes y la manzana: la mitad de la frescura. La pasta: el
cuarto de las féculas. El pescado y el yogur: el cuarto de los
constructores. El agua: el vaso. Están servidas todas las partes del
plato.
\end{solution}

\begin{solution}{exo:g3:balanced-diet:6}
El cuerpo lleva toda la noche funcionando sin entregas; el desayuno lo
reabastece. Sin él, a media mañana la \omterm{prop:g3:balanced-diet:jobs}{energía} se queda corta: tripas
que suenan, atención que se va, poca fuerza para jugar.
\end{solution}

\begin{solution}{exo:g3:balanced-diet:7}
El niño --- el cuerpo de un niño se sigue construyendo más alto y más
grande, además de repararse y de correr, así que el material de
construcción le importa proporcionalmente más.
\end{solution}

\begin{solution}{exo:g3:balanced-diet:8}
Respuesta abierta. La familia que más se olvida es la de la fruta y la
verdura; la que aparece de más, los alimentos dulces.
\end{solution}

\begin{solution}{exo:g3:balanced-diet:9}
Ningún alimento está prohibido: el equilibrio se juzga en días, no en
bocados. La regla es que los caprichos tienen que quedarse pequeños y
no ocupar el sitio de las familias --- la pregunta es ``¿qué familia
se me ha olvidado?'', no ``¿qué alimento es malo?''.
\end{solution}

\begin{solution}{exo:g3:balanced-diet:10}
El deporte gasta \omterm{prop:g3:balanced-diet:jobs}{energía} de más en el movimiento, y el frío gasta
\omterm{prop:g3:balanced-diet:jobs}{energía} de más en mantener caliente el cuerpo. Al gastar más, el
cuerpo pide más: el hambre crece con los gastos del día.
\end{solution}
